O Exame Nacional do Ensino Médio (ENEM) é a principal forma de entrada para o
ensino superior no Brasil, porém os resultados mais recentes revelam uma queda
persistente no desempenho em Matemática. Avaliações internacionais, como o PISA
2022, confirmam esse cenário crítico ao evidenciarem que muitos estudantes
brasileiros apresentam dificuldades em Matemática, e apenas uma pequena parcela
demonstra elevado nível de proficiência na disciplina. Pesquisas indicam que
essas dificuldades decorrem não apenas de lacunas de conhecimento, mas também de
estratégias cognitivas inadequadas, baixa compreensão dos enunciados e
desengajamento dos alunos. Apesar da ênfase da Base Nacional Comum Curricular (BNCC) no
Pensamento Computacional (PC) como habilidade fundamental para a resolução de
problemas, poucas ferramentas educacionais integram efetivamente o PC a
estratégias pedagógicas engajadoras. Este trabalho apresenta o Cognitivo, um
software gamificado fundamentado nos quatro pilares do PC (decomposição,
abstração, reconhecimento de padrões e algoritmos) projetado para possibilitar o
aumento do engajamento dos estudantes em preparação para a prova de matemática
do ENEM. Para o desenvolvimento da solução, adotou-se a \textit{Design-Based
Research} (DBR), organizada em ciclos iterativos de identificação do problema,
desenvolvimento do artefato, avaliação e refinamento. A validação foi realizada
com professores de Matemática por meio de um grupo focal, e a avaliação com
estudantes foi estruturada com o Método de Amostragem de Experiência (ESM) e a
Escala de Engajamento do Usuário (UES). Os resultados da validação com os
professores indicaram que o Cognitivo é viável e aplicável ao contexto
educacional, evidenciando o erro como elemento de reforço da aprendizagem, as
dicas como suporte ao esforço cognitivo e as mecânicas de gamificação como
promotoras da autorreflexão.

\palavraschave{Pensamento Computacional; Gamificação; Engajamento; Resolução de Problemas; ENEM; Matemática.}